\section{Proposed extension interfaces}

The following interfaces are design schemas. They name the obligations a
production implementation must expose; they are not compiled definitions from
the microprototype.

\begin{Code}
RuleProposal := {
  parentRequest, parentBranch, requiredContext,
  expressionOrDeclarationPlan,
  newObligations, writableMetavariables,
  applicabilityEvidence?, heuristicCost,
  continuation
}

DomainReply :=
  | Unsupported(reason)
  | RankingHint(score, observations)
  | NecessaryConstraint(formula, implicationProof)
  | Counterexample(input, validityProof, violationProof)
  | NoCompletion(certificate, exactScope)

ProofReply :=
  | Proof(term, assignmentDelta)
  | Refutation(term, assignmentDelta)
  | ReducedGoal(goal, reconstructionProof)
  | Unknown(reason, consumedBudget)

AcceptedCandidate := opaque {
  frozenRequest,
  closedProgram, closedContractProof, checkedAuxiliaries,
  dependencyAudit, sourceReplay,
  immutableReceipt
}
\end{Code}

A proposal's claimed context or assignment delta is validated against the
parent branch. \code{NoCompletion} needs a certificate for its exact completion
scope. \code{RankingHint} has no logical authority. The accepted-result
constructor is private and reachable only through the acceptance transaction.
These boundaries let a fast unverified proposer coexist with a small checker.

\Needspace{17\baselineskip}
\section{The central microprototype mechanism}

The essential control structure of the tested \code{search} function is shown
below in pseudocode. The complete actual Lean source is included separately.

\begin{Code}
search(root, pending, fuel):
  if pending is empty:
      succeed only if root has no unresolved expression metavariables
  discard pending goals already assigned by earlier constraints
  fail this bounded alternative if fuel is exhausted
  choose the next pending goal in its own local context
  for alternative in [locals, intro, constructors, providers, local spines]:
      save the full MetaM state
      apply the alternative, obtaining new goals
      if search(root, newGoals ++ remainingGoals, fuel - 1):
          return success without undoing this branch
      restore the saved MetaM state
  return bounded failure
\end{Code}

The recursive call is inside the saved alternative. Moving it outside and
solving new goals independently would destroy contract-sensitive backtracking.
The production design adds explicit cursors, precise error categories,
dependency-directed ordering, universe and delayed-obligation checks, and
separate final acceptance. The example's terminal expression-metavariable test
alone is not that production acceptance boundary.

\section{Reproducing the artifacts}\label{app:reproduce}

The archive is organized as follows:
\begin{Code}
Leant2/
  README.md
  lean-toolchain
  article/Leant2.tex
  article/Leant2.pdf
  article/sections/*.tex
  article/references.tex
  prototype/MicroSynth.lean
  prototype/Certificates.lean
  experiments/cegis.py
  experiments/results.json
  receipts/*.submitted.lean
  receipts/check-summary.json
  receipts/README.md
  tools/check_remote.py
  tools/remote_check.wls
\end{Code}

With Lean 4.34.0 installed, the two files can be checked independently using:
\begin{Code}
lean prototype/MicroSynth.lean
lean prototype/Certificates.lean
\end{Code}
They import \code{Lean}; they do not require this proposed Leant2 package to
exist. In this work they were checked remotely, not by these local commands.
The included remote helpers submit the selected file without replacing its
imports, retain the response, and check its status. The summaries in
\code{receipts} are transcriptions of tool responses, not signed independent
certificates.

The finite experiment uses only Python's standard library:
\begin{Code}
python experiments/cegis.py --output experiments/results.json
\end{Code}
Rebuilding the article requires XeLaTeX, the listed LaTeX packages, Latin Modern Roman/Sans and DejaVu Sans Mono installed on the build system:
\begin{Code}
cd article
latexmk -xelatex -interaction=nonstopmode -halt-on-error Leant2.tex
\end{Code}
Font files are not included in the archive.

\section{Evidence ledger}

The repository snapshot for Leant is
\code{6bf05ad78c467989e68290f2d08bbed40802d485}; for the independently inspected
Djex repository it is
\code{e8778f4ebd63e1f9b9b410fa4de8d14a8a04c9e5}. Leant's README describes its
vendored Djex pin separately; the two HEAD revisions should not be confused
with a claim that Leant builds against the independently inspected Djex HEAD.

The expanded microprototype check request was
\code{1490aa69-0708-478a-903d-b85ce0adbf78}. The fresh certificate check was
\code{a4f8d54b-9c4d-42a1-96cd-a6186919d756}; the cached repeat was
\code{e80c576f-ad6b-406d-abf7-f4b41fa6af5c}. The response reports
\code{lean-4.34.0}. Full executor identifiers, outcomes, warnings, and the
initial syntax failure are recorded in the JSON summary.

The claims in this article have four different evidence levels. Repository
observations are tied to the cited snapshots. Lean language/API descriptions
are tied to the cited documentation or versioned source. Experimental results
are tied to the included scripts, submitted source, and check summaries.
Production architecture, coverage priorities, and performance expectations are
proposals, not experimental measurements. The separation is deliberate: it
makes the document usable as an implementation specification without mistaking
its planned modules for delivered code.
